\ifdefined\iscombined
	\lecturetitle{System Plans, Unit Plans, and Root Cause Analyses}
\else
\documentclass{article}
\usepackage{changepage}
\usepackage{centernot}
\usepackage{listings}
\usepackage{amsthm}
\usepackage{braket}
\usepackage{comment}
\usepackage{algorithm}
\usepackage{algpseudocode}
\usepackage{amsmath}
\usepackage{amsfonts}
\usepackage{sectsty}
\usepackage{hyperref}
\usepackage{fancyhdr}
\usepackage[top=1in,bottom=1in,left=1in,right=1in]{geometry}
\usepackage{float}
\usepackage{graphicx}
\usepackage{tabularx}
\usepackage{mathtools}
\usepackage{tikz}
\usetikzlibrary{backgrounds, calc, positioning}

\date{
	\vspace{-.75em}
	{\large \today}
}

\title{
	\vspace*{-1.5em}
	{\Huge Lecture 27} \\
	\vspace{-1em}
	\rule{\textwidth/4}{.01em} \\
	\vspace{-.25em}
	{\Large System Plans, Unit Plans, and Root Cause Analyses}
	\vspace{-.75em}
}

\author{
	{\large Matthew Farah}
}

\hypersetup{
	colorlinks=true,
	linkcolor=blue,
	filecolor=magenta,
	urlcolor=blue,
	citecolor=blue,
	anchorcolor=blue
}

\fancyhead{}
\fancyhead[L]{\today}
\fancyhead[R]{Lecture 27 - System Plans, Unit Plans, and Root Cause Analyses}

\setlength{\parindent}{0cm}

\lstset{
	basicstyle=\ttfamily\small,
	frame=single,
	tabsize=4,
	breaklines=true,
	numbers=left,
	numberstyle=\tiny,
	numbersep=8pt,
	xleftmargin=1.5em,
	framexleftmargin=1.5em
}

\newcounter{example}
\renewcommand{\theexample}{\arabic{example}}

\newenvironment{example}[1][]{
	\refstepcounter{example}
	\begin{adjustwidth}{1.5em}{}
	\noindent\textbf{\ifx&#1&Example \theexample:\else Example \theexample \ (#1):\fi}
	\ignorespaces
}{
	\vspace{.5em}
	\end{adjustwidth}
}

\newcounter{subexample}
\renewcommand{\thesubexample}{\alph{subexample}}

\newenvironment{subexample}{
	\setcounter{step}{0}
	\refstepcounter{subexample}
	\vspace{1em}\par\noindent
	\begin{adjustwidth}{1.5em}{}
	\noindent\textbf{(\thesubexample)}
	\ignorespaces
}{
	\par
	\vspace{1em}
	\end{adjustwidth}
}

\newenvironment{solution}{
	\vspace{.5em}
	\par
	\begin{adjustwidth}{1.5em}{}
	\textbf{{Solution:}}
	\par
	\vspace{.5em}
}{
	\vspace{1em}
	\end{adjustwidth}
}

\newcounter{step}
\renewcommand{\thestep}{\Roman{step}}

\newenvironment{step}{
	\refstepcounter{step}
	\vspace{1em}\par\noindent
	\ignorespaces\noindent\textbf{(\thestep)}
	\ignorespaces
}{
	\par
	\vspace{1em}
}

\sectionfont{\fontsize{12}{12}\bfseries}
\subsectionfont{\fontsize{10}{12}\normalfont}
\subsubsectionfont{\fontsize{10}{12}\normalfont}

\begin{document}

\maketitle
\pagestyle{fancy}

\fi


\begin{itemize}
	\item \emph{Aside: Root cause analysis question likely on exam}.
	\begin{itemize}
		\item \emph{Aside: Read M7.L4.S/UT Plan and QATeam lecture slides for more details on fish-bone diagrams and components of system/unit plans, don't wanna type it out here}. \\
	\end{itemize}
	\item A testing plan is not the same as a testing schedule.
	\begin{itemize}
		\item A given testing plan may have multiple valid testing schedules.
	\end{itemize}
	\item A testing plan is either a:
	\begin{enumerate}
		\item System plan.
		\begin{itemize}
			\item Plans for testing the entire system as a whole.
		\end{itemize}
		\item Unit plan.
		\begin{itemize}
			\item Plans for testing an individual component.
			\item Essentially a scaled-down system test.
		\end{itemize}
	\end{enumerate}
	%TODO: Maybe elaborate on the components.
	\item System plans are comprised of four main components.
	\begin{enumerate}
		\item General Information.
		\item The plan.
		\item Specifications and Evaluations.
		\item Test Descriptions.
	\end{enumerate}
	%TODO: Also elaborate here.
	\item Unit plans are also comprised of four main components.
	\begin{enumerate}
		\item The plan.
		\item Business and structural function testing.
		\item An interface test description.
		\item Test progression.
	\end{enumerate}
	\item Key takeaways from previous testing plans can be analyzed using the 8D approach.
	\item A root cause analysis should be done when a bug has been caught.
	\begin{itemize}
		\item Done to help reduce number of bugs in current/future development.
	\end{itemize}
	\item Root cause analyses seek to answer the following five core questions regarding a:
	\begin{enumerate}
		\item What happened?
		\item When did it happen?
		\begin{itemize}
			\item Includes both when it was first introduced and when it was first identified.
		\end{itemize}
		\item Why did it happen?
		\item How can a similar bug be prevented from being introduced in the future?
	\end{enumerate}
	\item Root cause analyses can be modelled using fish-bone diagrams.
	%TODO: Elaborate
	\item Root cause analysis:
	\begin{itemize}
		\item Begins with fault selection.
		\item Groups different faults according to the severity.
	\end{itemize}
\end{itemize}

\ifdefined\iscombined
	\newpage
\else
	\end{document}
\fi
